\section{Quasi-invariant polynomials for $S_n$}\label{sec:2}

We start with the definition of quasi-invariant polynomials. Let $\k$ be a field, and let $S_n$ act on $\k[x_1,\dots,x_n]$ by permuting the variables. Let $s_{i,j}\in S_n$ denote the transposition swapping the $i$-th and $j$-th indices.

\begin{definition}
    Let $n$ be a positive integer. A polynomial $P\in\k[x_1,\dots,x_n]$ is $m$-quasi-invariant if for all transpositions $s_{i,j}\in S_n$, $(1-s_{i,j})P$ is divisible by $(x_i-x_j)^{2m+1}$. We denote by $Q_m(n,\k)$ the space of all $m$-quasi-invariants in $\k[x_1,\dots,x_n]$. 
\end{definition}

Let $\k[x_1,\dots,x_n]^{S_n}$ denote the space of polynomials in $\k[x_1,\dots,x_n]$ which are invariant with respect to the action of $S_n$. Some useful properties of $Q_m(n,\k)$ are described in the following proposition. 

\begin{proposition}[\cite{etingof2002lectures}]\label{prop:2.2}
    \phantom{.}

    \begin{enumerate}
        \item $\k[x_1,\dots,x_n]^{S_n}\subset Q_m(n,\k)$, $Q_0(n,\k)=\k[x_1,\dots,x_n]$, $Q_m(n,\k)\subset Q_{m'}(n,\k)$ if $m>m'$.
        \item $Q_m(n,\k)$ is a ring.
        \item $Q_m(n,\k)$ is a finitely generated module over $\k[x_1,\dots,x_n]^{S_n}$.
    \end{enumerate}
\end{proposition}

Note that $Q_m(n,\k)$ admits a grading by the degrees of the polynomials. Thus we may define its Hilbert series.

\begin{definition}
    The Hilbert series of $Q_m(n,\k)$ is
    $$H_m(t)=\sum_{d\geq 0}\dim Q_m(n,\k)[d]t^d$$
    where $Q_m(n,\k)[d]$ is the graded component of $Q_m(n,\k)$ consisting of polynomials of homogeneous degree $d$.
\end{definition}

Note that by the fundamental theorem of symmetric polynomials, $\k[x_1,\dots,x_n]^{S_n}$ is a free polynomial algebra generated by symmetric polynomials $P_1,\dots,P_n$ of degrees $1,\dots,n$, respectively \cite{humphreys_1990}. Thus since $Q_m(n,\k)$ is a finitely generated module over $\k[x_1,\dots,x_n]^{S_n}=\k[P_1,\dots,P_n]$ by Proposition \ref{prop:2.2}, we can write
$$H_m(t)=\frac{G_m(t)}{\prod_{d=1}^n(1-t^d)}$$
where $G_m(t)$ is a polynomial by Hilbert's syzygy theorem. We say that $G_m(t)$ is the Hilbert polynomial associated with $H_m(t)$. Note that describing the Hilbert polynomial is related to describing the generators and relations of $Q_m(n,\k)$ as a $\k[x_1,\dots,x_n]^{S_n}$-module.

From now on, we will assume the characteristic of $\k$ does not divide $n!$, so that the representation theory of $S_n$ over $\k$ is non-modular. The following lemma and proposition are standard results in the study of quasi-invariants, for example they a direct consequence of the decomposition of $Q_m(n,\C)$ as a representation of the spherical rational Cherednik algebra (see \cite{berest2003cherednik}) in the case $\k=\C$. 

\begin{lemma}\label{lem:2.4}
    $Q_m(n,\k)$ is a representation of $S_n$.
\end{lemma}

Now, consider a graded component $Q_m(n,\k)[d]$ of the quasi-invariants. It is a subrepresentation of $Q_m(n,\k)$, and moreover it is finite dimensional since it is a subspace of the finite dimensional space $\k[x_1,\dots,x_n][d]$. Thus, since we are working in the non-modular case, by Maschke's theorem it decomposes as a direct sum
$$Q_m(n,\k)[d]=\bigoplus_{\tau\in\Irrep(S_n)}Q_m(n,\k)[d]_\tau$$
where $\Irrep(S_n)$ is the set of irreducible representations of $S_n$ and $Q_m(n,\k)[d]_\tau$ is the subspace of $Q_m(n,\k)[d]$ on which $S_n$ acts by a direct sum of copies of $\tau$. Then for $\tau\in\Irrep(S_n)$, let us define
$$Q_m(n,\k)_\tau=\bigoplus_{d\geq 0}Q_m(n,\k)[d]_\tau$$
to be the isotypic components of $Q_m(n,\k)$.

\begin{proposition}\label{prop:2.5}
\begin{enumerate}
    \item $Q_m(n,\k)_\tau$ is a module over $\k[x_1,\dots,x_n]^{S_n}$.\label{prop:2.51}
    \item We have the decomposition
    $$Q_m(n,\k)=\bigoplus_{\tau\in\Irrep(S_n)}Q_m(n,\k)_\tau$$
    as a $\k$-vector space and hence as a $\k[x_1,\dots,x_n]^{S_n}$-module as well.\label{prop:2.52}
\end{enumerate}
\end{proposition}

By the above proposition, to study the generators and relations of $Q_m(n,\k)$ it suffices to study the generators and relations of $Q_m(n,\k)_\tau$. This will be our primary method of studying $Q_m(n,\k)$ throughout the rest of this paper. Let us denote by triv and sign the trivial and sign representations of $S_n$, respectively. We deal with these two cases here.

\begin{proposition}\label{prop:2.6} Assume $\char\k\neq 2$. As $\k[x_1,\dots,x_n]^{S_n}$-modules, we have
\begin{enumerate}
    \item $Q_m(n,\k)_\triv$ is freely generated by $1$.\label{prop:2.61}
    \item $Q_m(n,\k)_\sgn$ is freely generated by $\prod_{i<j}(x_i-x_j)^{2m+1}$.\label{prop:2.62}
\end{enumerate}
\end{proposition}
\begin{proof}
    \ref{prop:2.61}~Note that the statement that $S_n$ acts by the trivial representation on $P$ is equivalent to the statement that $P$ is $S_n$-invariant, so $Q_m(n,\k)_\triv\subset\k[x_1,\dots,x_n]^{S_n}$. On the other hand, for any $P\in\k[x_1,\dots,x_n]^{S_n}$ and any $s_{i,j}\in S_n$ we have $(1-s_{i,j})P=P-P=0$ is divisible by $(x_i-x_j)^{2m+1}$, so $Q_m(n,\k)_\triv=\k[x_1,\dots,x_n]^{S_n}$.
    
    \ref{prop:2.62}~Let $P\in Q_m(n,\k)_\sgn$ and let $s_{i,j}\in S_n$. Then since $P$ is in the sign representation we have $s_{i,j}P=-P$, so $(1-s_{i,j})P=2P$ is divisible by $(x_i-x_j)^{2m+1}$. So $P$ is divisible by $(x_i-x_j)^{2m+1}$ for all $i$ and $j$, and thus it is divisible by $\prod_{i<j}(x_i-x_j)^{2m+1}$. 
    
    Let us write $P=K\prod_{i<j}(x_i-x_j)^{2m+1}$ for some polynomial $K$. Note that 
    $$s_{i',j'}\prod_{i<j}(x_i-x_j)^{2m+1}=-\prod_{i<j}(x_i-x_j)^{2m+1}$$
    for all $s_{i',j'}\in S_n$, so we have 
    $$s_{i',j'}P=-P=-K\prod_{i<j}(x_i-x_j)^{2m+1}=s_{i',j'}K\bigg(s_{i',j'}\prod_{i<j}(x_i-x_j)^{2m+1}\bigg)=-s_{i',j'}K\prod_{i<j}(x_i-x_j)^{2m+1}$$
    so $s_{i',j'}K=K$. Since $S_n$ is generated by transpositions, $K$ is $S_n$-invariant, and thus $P$ is in the $\k[x_1,\dots,x_n]^{S_n}$-module generated by $\prod_{i<j}(x_i-x_j)^{2m+1}$. Conversely, for any $K\in\k[x_1,\dots,x_n]^{S_n}$ and $s_{i',j'}\in S_n$, we have 
    $$s_{i',j'}\Big(K\prod_{i<j}(x_i-x_j)^{2m+1}\Big)=-K\prod_{i<j}(x_i-x_j)^{2m+1}$$
    and $$(1-s_{i',j'})\Big(K\prod_{i<j}(x_i-x_j)^{2m+1}\Big)=2K\prod_{i<j}(x_i-x_j)^{2m+1}$$ is divisible by $(x_{i'}-x_{j'})^{2m+1}$, so the $\k[x_1,\dots,x_n]^{S_n}$-module generated by $\prod_{i<j}(x_i-x_j)^{2m+1}$ is contained in $Q_m(n,\k)_\sgn$. So $Q_m(n,\k)_\sgn$ is exactly this module. Freeness then follows from the fact that $\k[x_1,\dots,x_n]$ is an integral domain.
\end{proof}

\begin{remark}\label{rem:2.7}
    It was proved in \cite{Ren_2020} that the Hilbert series of $Q_m(2,\k)$ is $\frac{1+t^{2m+1}}{(1-t)(1-t^2)}$, regardless of the characteristic of $\k$. Proposition \ref{prop:2.6} provides an alternate proof of this result in the case $\char \k\neq 2$. Namely, note that the only irreducible representations of $S_2$ are triv and sign. Thus the generators of $Q_m(2,\k)$ are 1 (of degree 0) and $(x_1-x_2)^{2m+1}$ (of degree $2m+1$). Since these generators belong to different irreducible representations of $S_2$, there are no relations between them, and thus the Hilbert polynomial is $1+t^{2m+1}$ and the Hilbert series is $\frac{1+t^{2m+1}}{(1-t)(1-t^2)}$.
\end{remark}

\section{The Hilbert series of $Q_m(3,\k)$}\label{sec:3}

The Hilbert series of $Q_m(3,\C)$ was computed in \cite{felder2003action} and is as follows:
\begin{equation}\label{eq:3.1}
    H_m(t)=\frac{1+2t^{3m+1}+2t^{3m+2}+t^{6m+3}}{(1-t)(1-t^2)(1-t^3)}.\tag{3.1}
\end{equation}

In this section, we work toward an explicit description of the Hilbert series of $Q_m(3,\k)$ in characteristic $p>3$. Note that there are three irreducible representations of $S_3$: triv, sign, and the 2-dimensional standard representation which we will denote by std. Proposition \ref{prop:2.6} computes the generators of $Q_m(3,\k)_\triv$ and $Q_m(3,\k)_\sgn$, so to compute the Hilbert series of $Q_m(3,\k)$ we only need to consider the space $Q_m(3,\k)_\std$. Both this section and the next section will be largely focused on understanding this space.

This section is dedicated to proving the main result of this paper, which is the following theorem, conjectured in \cite{Ren_2020}. Note the similarity to (\ref{eq:3.1}).

\begin{theorem}\label{thm:3.1}
    Let $\k$ be a field of characteristic $p>3$. Then $Q_m(3,\k)_\std$ is freely generated by standard representations of degree $d$ and $6m+3-d$ where $d$ is an integer depending on $m,p$ that satisfies $2m+1\leq d\leq 3m+1$. In particular, the Hilbert series of $Q_m(3,\k)$ is
    $$H_m(t)=\frac{1+2t^d+2t^{6m+3-d}+t^{6m+3}}{(1-t)(1-t^2)(1-t^3)}.$$
    
\end{theorem}

To start, let $V$ be a copy of the standard representation, and let $s_{i,j}\in S_3$. Then $V$ contains a one dimensional $-1$ eigenspace of $s_{i,j}$. Let us denote this eigenspace by $V_{i,j}^-$. Then we have

\begin{lemma}\label{lem:3.2}
    Let $V\subset Q_m(3,\k)_\std$ be a copy of the standard representation, and let $P\in V_{i,j}^-$. Then we have $P+sP+s^2P$=0 where $s=(1\,2\,3)\in S_3$ and $P=(x_i-x_j)^{2m+1}K$ for some polynomial $K$ that is invariant under the action of $s_{i,j}$. Conversely, let $K$ be an $s_{1,2}$-invariant polynomial such that 
    $$(x_1-x_2)^{2m+1}K+(x_2-x_3)^{2m+1}sK+(x_3-x_1)^{2m+1}s^2K=0.$$
    Then $(x_1-x_2)^{2m+1}K$ either belongs to $Q_m(3,\k)_\sgn$ or the $-1$ eigenspace of $s_{1,2}$ in some copy of std inside $Q_m(3,\k)_\std$.
\end{lemma}
\begin{proof}
    The proof of the first statement is analogous to the proof of Proposition \ref{prop:2.6}. Since $P\in V_{i,j}^-$, we have $(1-s_{i,j})P=2P$ is divisible by $(x_i-x_j)^{2m+1}$, so $P=(x_i-x_j)^{2m+1}K$ for some polynomial $K$. Then since $P$ and $(x_i-x_j)^{2m+1}$ are both $s_{i,j}$-antiinvariant, $K$ must be $s_{i,j}$-invariant. The fact that $P+sP+s^2P=0$ is a direct consequence of $P$ being an element of the standard representation.
    
    For the second statement, note that since $K$ is $s_{1,2}$-invariant, $sK=s_{1,3}K$ and $s^2K=s_{2,3}K$. So we have
    $$(1-s_{1,3})\Big((x_1-x_2)^{2m+1}K\Big)=(x_1-x_2)^{2m+1}K-(x_3-x_2)^{2m+1}sK=-(x_3-x_1)^{2m+1}s^2K,$$
    $$(1-s_{2,3})\Big((x_1-x_2)^{2m+1}K\Big)=(x_1-x_2)^{2m+1}K-(x_1-x_3)^{2m+1}s^2K=-(x_2-x_3)^{2m+1}sK$$
    so $(x_1-x_2)^{2m+1}K$ is $m$-quasi-invariant. On the other hand, since $s_{1,2}$ sends the vector space spanned by $(x_1-x_2)^{2m+1}K$ to itself and $s_{1,3}$ and $s_{2,3}$ send $(x_1-x_2)^{2m+1}K$ to polynomials in the space spanned by $(x_1-x_2)^{2m+1}K$ and $(x_2-x_3)^{2m+1}sK$, the span of the $S_3$ orbit of $(x_1-x_2)^{2m+1}K$ is at most 2 dimensional. If it is 1 dimensional, then $(x_1-x_2)^{2m+1}K$ is contained in a copy of sign inside $Q_m(3,\k)_\sgn$ since $s_{1,2}$ acts by $-1$. If it is 2 dimensional, then since the only 2 dimensional irreducible representation of $S_3$ is std, $(x_1-x_2)^{2m+1}K$ belongs to some copy of std inside $Q_m(3,\k)_\std$. That it is in the $-1$ eigenspace of $s_{1,2}$ is clear by applying $s_{1,2}$ to $(x_1-x_2)^{2m+1}K$.
\end{proof}

Note that since $Q_m(3,\k)_\std$ is a graded $\k[x_1,x_2,x_3]^{S_3}$-module, we may assume its generators are homogeneous. We may also describe generating representations of $Q_m(3,\k)_\std$, rather than generating elements. 

\begin{corollary}\label{cor:3.3}
    Let $V$ be a generating representation of $Q_m(3,\k)_\std$ and let $P\in V_{i,j}^-$. Let us write $P=(x_i-x_j)^{2m+1}K$. Then $K$ is not divisible by any nonconstant symmetric polynomial.
\end{corollary}
\begin{proof}
    Since $P$ is in $Q_m(3,\k)_\std$, it satisfies $P+sP+s^2P=0$ where $s=(1\,2\,3)$. Assume for contradiction that $K$ is divisible by some symmetric polynomial $Q$. Then it is clear that \linebreak
    $P/Q+s(P/Q)+s^2(P/Q)=0$ and $P/Q=(x_i-x_j)^{2m+1}(K/Q)$, so $P/Q$ is contained in $Q_m(3,\k)$ by Lemma \ref{lem:3.2}. But then $P$ is in the $\k[x_1,x_2,x_3]^{S_3}$-module generated by $P/Q$, so it is generated by an element of lower degree and thus is not contained in a generating representation.
\end{proof}

For a standard representation $V$ consisting of homogeneous polynomials of degree $d$, let us say its degree is $\deg V=d$. The following lemma makes use of this, and will prove to be useful both in the proof of Theorem \ref{thm:3.1} and in Section \ref{sec:4}.

\begin{lemma}\label{lem:3.4}
    Let $V,W$ be distinct generating representations of $Q_m(3,\k)_\std$. Let $A\in V_{1,2}^-,B\in W_{1,2}^-$. Then $As_{2,3}B-Bs_{2,3}A$ is a nonzero element of $Q_m(3,\k)_\sgn$ and we have $\deg V+\deg W\geq 6m+3$.
\end{lemma}
\begin{proof}
    $As_{2,3}B-Bs_{2,3}A$ is contained in $Q_m(3,\k)$ because the quasi-invariants form a ring by Proposition \ref{prop:2.2}. It is also contained in a sign representation since it spans the representation $\wedge^2\std=\sgn$ inside $V\otimes W$. Thus it is contained in $Q_m(3,\k)_\sgn$, so it remains to show that it is nonzero.
    
    By Lemma \ref{lem:3.2} we can write $A=(x_1-x_2)^{2m+1}K,B=(x_1-x_2)^{2m+1}L$ for $s_{1,2}$-invariant polynomials $K$ and $L$. Then we have
    \begin{align*}
        As_{2,3}B-Bs_{2,3}A&=(x_1-x_2)^{2m+1}K(x_1-x_3)^{2m+1}s_{2,3}L-(x_1-x_2)^{2m+1}L(x_1-x_3)^{2m+1}s_{2,3}K\\
        &=(x_1-x_2)^{2m+1}(x_1-x_3)^{2m+1}(Ks_{2,3}L-Ls_{2,3}K).
    \end{align*}
    So $As_{2,3}B-Bs_{2,3}A$ is zero if and only if $Ks_{2,3}L=Ls_{2,3}K$, or in other words if $Ks_{2,3}L$ is $s_{2,3}$-invariant. Let us assume that $Ks_{2,3}L$ is $s_{2,3}$-invariant. We will show that this implies $K=L$ up to a scalar, which contradicts $V$ and $W$ being distinct.
    
    Decompose $K$ into its irreducible factors, and let $Q$ be such a factor. We will show that $Q$ is also a factor of $L$. Since $Ks_{2,3}L$ is $s_{2,3}$-invariant, $s_{2,3}Q$ must be a factor of either $K$ or $s_{2,3}L$. If it is a factor of $s_{2,3}L$, then $P$ is a factor of $L$, as desired. Otherwise, $s_{2,3}Q$ is a factor of $K$. Then we also have that $s_{1,2}Q$ and $s_{1,2}s_{2,3}Q$ are factors of $K$ since $K$ is $s_{1,2}$-invariant. Then $s_{2,3}s_{1,2}Q$ is a factor of $Ks_{2,3}L$. If it is a factor in $s_{2,3}L$, then note that $s_{1,3}s_{2,3}s_{1,2}Q$ is a factor of $s_{2,3}L$ since $L$ being $s_{1,2}$-invariant implies $s_{2,3}L$ is $s_{1,3}$-invariant. But we have $s_{1,3}s_{2,3}s_{1,2}=s_{2,3}$, so $s_{2,3}Q$ is a factor of $s_{2,3}L$, which implies $Q$ is a factor of $L$, as desired. So we may assume $s_{2,3}s_{1,2}Q$ is a factor of $K$. Then we also have $s_{1,2}s_{2,3}s_{1,2}Q=s_{1,3}Q$ is a factor of $K$ since $K$ is $s_{1,2}$-invariant, and thus $K$ is divisible by $sQ$ for all $s\in S_3$. 
    
    So $K$ is divisible by the least common multiple of all of the $sQ$, which we will denote by $S$. $S_3$ fixes $S$ as a factor of $K$, so it must act on $S$ by scalars. It cannot act on $S$ trivially, as then $S$ would be a symmetric polynomial factor of $K$, and $K$ has no symmetric polynomial factors by Corollary \ref{cor:3.3}. So $S_3$ must act on $S$ by sign. Note also that we can only have one such factor of $K$ on which $S_3$ acts by sign, as otherwise their product would be a symmetric polynomial factor of $K$. So every factor of $K$ is also a factor of $L$ except for possibly a factor in the sign representation. We may apply everything above to $L$ as well to get the same result in the other direction, yielding a bijection between all irreducible factors of $K$ and $L$ except maybe a factor of the sign representation in either case. Since $K$ and $L$ are both $s_{1,2}$-invariant, acting on both of them by $s_{1,2}$ makes it clear that either they both have an extra factor of the sign representation or they both do not.
    
    If $K$ and $L$ do not have an additional factor in the sign representation, then we have a bijection between all irreducible factors of $K$ and $L$, and thus they are equal up to a scalar, as desired. Otherwise, note that by Proposition \ref{prop:2.2}, the sign factors lie in $Q_0(3,\k)$, and by Proposition \ref{prop:2.6} they are symmetric polynomial multiples of $(x_1-x_2)(x_2-x_3)(x_3-x_1)$. But these symmetric polynomial multiples must be 1, as otherwise $K$ and $L$ would be divisible by nonconstant symmetric polynomials, which violates Corollary \ref{cor:3.3}. So the sign factors are both $(x_1-x_2)(x_2-x_3)(x_3-x_1)$. Thus we still have a bijection between the factors of $K$ and $L$, so $K$ and $L$ are equal up to a scalar, as desired.
    
    Note that $As_{2,3}B-Bs_{2,3}A$ has degree $\deg V+\deg W$. Since $Q_m(3,\k)_\sgn$ is generated by a polynomial of degree $6m+3$ by Proposition \ref{prop:2.6}, we have $\deg V+\deg W\geq 6m+3$.
\end{proof}

The above lemma brings us closer to the proof of Theorem \ref{thm:3.1}, however we still need a couple more lemmas.

\begin{lemma}\label{lem:3.5}
    Assume that there exist generating representations $V,W$ of $Q_m(3,\k)_\std$ such that $\deg V+\deg W=6m+3$. Then $Q_m(3,\k)_\std$ is a free module over $\k[x_1,x_2,x_3]^{S_3}$ generated by $V$ and $W$.
\end{lemma}
\begin{proof}
    Assume for contradiction that there exists some other generating representation $U$ of \linebreak
    $Q_m(3,\k)_\std$. By Lemma \ref{lem:3.4} we have $\deg U\geq \deg W$. Let $A\in V_{1,2}^-,B\in W_{1,2}^-,C\in U_{1,2}^-$. Then note that by Lemma \ref{lem:3.4} we have
    $$As_{2,3}B-Bs_{2,3}A=c\prod_{i<j}(x_i-x_j)^{2m+1}$$ 
    where $c$ is a nonzero constant, since it is a degree $6m+3$ polynomial in $Q_m(3,\k)_\sgn$. We also have
    $$As_{2,3}C-Cs_{2,3}A=Q\prod_{i<j}(x_i-x_j)^{2m+1}$$
    for some symmetric polynomial $Q$. But then instead of $U$, we may take the representation generated by $cC-QB$ as the generating representation of $Q_m(3,\k)_\std$. But we have
    \begin{align*}
        As_{2,3}(cC-QB)-(cC-QB)s_{2,3}A&=c(As_{2,3}C-Cs_{2,3}A)-Q(As_{2,3}B-Bs_{2,3}A)\\
        &=(cQ-Qc)\prod_{i<j}(x_i-x_j)^{2m+1}\\
        &=0
    \end{align*}
    which contradicts Lemma \ref{lem:3.4}. If $Q_m(3,\k)_\std$ was not a free module then there would be nonzero symmetric polynomials $R_1,R_2$ such that $R_1A=R_2B$. But then we would have 
    $$R_1As_{2,3}(R_2B)-R_2Bs_{2,3}(R_1A)=R_1As_{2,3}(R_1A)-R_1As_{2,3}(R_1A)=0$$
    and 
    $$R_1As_{2,3}(R_2B)-R_2Bs_{2,3}(R_1A)=R_1R_2(As_{2,3}B-Bs_{2,3}A)$$
    which cannot happen since none of $R_1,R_2,as_{2,3}b-bs_{2,3}a$ are zero.
\end{proof}
